\documentclass[11pt]{article}

\usepackage[T1]{fontenc}
\usepackage{lmodern}
\usepackage[margin=1.1in]{geometry}
\usepackage{booktabs}
\usepackage{array}
\usepackage{enumitem}
\usepackage{amsmath}
\usepackage{amssymb}
\usepackage[hidelinks]{hyperref}
\usepackage{microtype}
\usepackage{tabularx}

\newcommand{\sys}[1]{\textbf{#1}}
\newcommand{\code}[1]{\texttt{#1}}
\newcolumntype{L}{>{\raggedright\arraybackslash}X}

\title{Evaluating Two-Stream Symbolic Music Generation:\\ Experiment Plan Report}
\author{}
\date{\today}

\begin{document}
\maketitle

\noindent\emph{Companion document to \code{EXPERIMENTS.md} (the
operational protocol).}

\section{Setup}

\paragraph{Tasks.} Generation of symbolic music with two paired
streams:
\begin{enumerate}[nosep]
  \item melody and chord pair. TRAINING corpus: POP909
    ($\sim$900 songs) $+$ Nottingham ($\sim$1{,}000 two-track tunes)
    combined, which roughly doubles the melody/chord training data;
    EVALUATION stays POP909-only (\S1, held-out ids), so metrics remain
    interpretable in the target domain and Nottingham contributes
    training signal only;
  \item drum and non-drum pair (LA corpus; evaluated on held-out RWC).
\end{enumerate}

\paragraph{Generation modes.}
\begin{enumerate}[nosep]
  \item co-generation (co-gen);
  \item marginal generation (marg.\ gen);
  \item conditional generation (cond.\ gen).
\end{enumerate}

\paragraph{Baselines.}
\begin{center}
\begin{tabularx}{\textwidth}{@{}lL@{}}
\toprule
ID & Construction \\
\midrule
\sys{S0} & pretrained single-stream CP transformer, trained on the LA dataset \\
\sys{S1} & S0 finetuned on the COMBINED POP909$+$Nottingham melody--chord data (merged streams), i.e.\ the same corpus the duet systems train on \\
\sys{S-mel}\,/\,\sys{S-chord} & S0 finetuned on \emph{one} stream of the combined melody--chord data \\
\sys{Y} & YinYang: frozen S0 $+$ low-rank cross-attention adapters $+$ LoRA; conditional generation \\
\bottomrule
\end{tabularx}
\end{center}

\paragraph{Our models.}
\begin{center}
\begin{tabularx}{\textwidth}{@{}lL@{}}
\toprule
ID & Model \\
\midrule
\sys{A.1} & DuetAttn (\code{M2CIntraCrossAttn}) \\
\sys{A.2} & DuetBlockDiffusion (\code{M2CDuetBlockDiffusion}) \\
\sys{B.1} & DuetAnticipatory (\code{M2CDuetAnticipatory}) \\
\sys{C.1} & \code{M2CDuetRehearsal} \\
\sys{C.2} & \code{M2CDuetPrefix} \\
\bottomrule
\end{tabularx}
\end{center}

\section{Research questions}

\begin{description}[nosep]
  \item[RQ1 --- Co-generation.] Does explicit two-stream joint modeling
    beat a merged-stream model at generating both parts together?
  \item[RQ2 --- Marginals.] Does joint training \emph{retain} the
    ability to generate each stream alone (a non-inferiority question,
    not a superiority one)?
  \item[RQ3 --- Conditional generation.] Given the full partner
    stream, how does the \emph{conditioning horizon} (no future /
    bounded future / unbounded future) shape quality, and do
    fully-trained conditional duets beat parameter-efficient adapter
    conditioning?
  \item[RQ4 --- Mechanisms.] Which architectural components are
    causally load-bearing: the within-frame refinement mechanism, the
    training objective behind it, and the MoE routing?
\end{description}

\section{Experiment blocks}

Table~\ref{tab:training-requirements} lists, per experiment, every
model that must be trained (or already exists) and the dataset it is
trained on. Checkpoints are reused across experiments wherever the
same system appears; only the first row that introduces a checkpoint
carries a training cost.

\begin{table}[ht]
\centering
\small
\begin{tabularx}{\textwidth}{@{}llLl@{}}
\toprule
Exp. & System & Training data & Status \\
\midrule
E1 & A.2 (drum/non-drum) & LA drum $+$ non-drum pair (cp16) & trained (98k) \\
E1 & A.2 (melody/chord) & COMBINED melody $+$ chord pair (POP909 $+$ Nottingham, cp4) & trained \\
E1 & B.1 (drum/non-drum) & LA drum $+$ non-drum pair & trained \\
E1 & B.1 (melody/chord) & COMBINED melody $+$ chord pair, $k{=}16$ & trained \\
E1 & S0 & LA merged single stream (pretraining) & trained \\
E1 & S1 & COMBINED merged, program-tagged (cp16); finetune of S0 & trained (short recipe) \\
\midrule
E2 & A.2 & --- (reuses the E1 checkpoints) & --- \\
E2 & S-mel & COMBINED melody-only stream; finetune of S0, short recipe & \textbf{to train} \\
E2 & S-chord & COMBINED chord-only stream; finetune of S0, short recipe & \textbf{to train} \\
E2 & S-drum / S-nondrum & LA drum-only / non-drum-only; finetunes of S0 & gated on melchord E2 \\
\midrule
E3 & A.2, B.1 & --- (reuse the E1 checkpoints) & --- \\
E3 & C.1, C.2 & LA drum $+$ non-drum pair & trained \\
E3 & Y-dn & LA drum-subset, 31k songs (adapters $+$ LoRA on frozen S0) & existing paper ckpt \\
E3 & Y-mc & Nottingham, 1{,}020 songs (adapters $+$ LoRA on frozen S0) & existing paper ckpt \\
E3 & (optional domain match) & a POP909-matched YinYang finetune, or a Nottingham-only evaluation slice; Nottingham is now part of the duet TRAINING mix, so Y-mc is no longer the only system seeing folk data & optional \\
\midrule
E4a & A.1 (drum/non-drum) & LA drum $+$ non-drum pair & trained \\
E4a & A.1 (melody/chord) & COMBINED melody $+$ chord pair & gated on drum/non-drum E4a \\
E4b & --- & decode-time only; reuses the A.2-melchord checkpoint & --- \\
E4c & A.2-dense & COMBINED melody $+$ chord pair (identical recipe to A.2-melchord) & \textbf{to train} \\
\midrule
E5 & B.1 at $k \in \{8, 32\}$ & LA drum $+$ non-drum pair ($k{=}16$ exists) & gated on E3 \\
\bottomrule
\end{tabularx}
\caption{Training requirements by experiment. Bold = new training
runs currently owed; ``gated'' = trained only if the triggering result
warrants it.}
\label{tab:training-requirements}
\end{table}

\subsection{E1 --- Co-generation (headline)}

Both streams are generated jointly from a 4-bar prompt of both.
Corpus-matched pairing: A.2-drumnondrum vs.\ \sys{S0} (both
LA-trained); A.2-melchord vs.\ \sys{S1} (both trained on the combined
POP909$+$Nottingham corpus) --- so
the measured gap is attributable to architecture alone. B.1
participates per-stream as the competing duet strategy.

Four pre-registered hypotheses map each burden the single-stream model
carries to a specific observable
(Table~\ref{tab:e1-hypotheses}).

\begin{table}[ht]
\centering
\small
\begin{tabularx}{\textwidth}{@{}lp{1.7cm}Lp{3.9cm}@{}}
\toprule
\# & Judges & Prediction & Primary metric \\
\midrule
H1 & each stream alone & consistent role \& texture across 24 bars
  (melody monophonic, chords block-voiced, stable density, no stream
  death); S* shows role leakage, density drift, or stream death &
  \code{survival\_min} \\
H2 & mutual fit & chords consonant with the simultaneous melody;
  non-drum onsets lock to drums &
  \code{chord\_tone\_cov\_delta} / \code{onset\_sync\_delta} \\
H3 & grammar & stream-appropriate statistics (harmonic rhythm
  $\approx$ 1 change/bar, contour smoothness); MoE experts separate by
  stream & \code{harmonic\_rhythm\_jsd} / \code{onset\_grid\_jsd} \\
H4 & listeners & an \emph{ordering} of listening margins:
  coherence $>$ musicality $\ge 0 \approx$ structure --- a uniform
  profile would indicate generic quality, not the mechanism &
  pairwise A/B per axis \\
\bottomrule
\end{tabularx}
\caption{E1 pre-registered hypotheses.}
\label{tab:e1-hypotheses}
\end{table}

Pre-registered counter-expectations: ties on local plausibility and
pitch-class histograms (the matched baseline saw the same data);
possible loss on the structure axis (A.2's known repetition tendency).

\subsection{E2 --- Marginal generation (guard; retention, not superiority)}

Claim: joint training did not cost the model its marginals. Opponent:
the per-stream \emph{specialist} --- same pretrained initialization,
same finetune recipe, same songs, single-modality data --- the
strictest fair opponent (each devotes full capacity to one marginal).
Success criterion: non-inferiority within a pre-registered margin
($\delta = 0.05$ on primary JSD grammar metrics). Outside-margin
results are reported as the \emph{quantified marginal cost} of joint
training; a better-than-specialist result is reported as positive
cross-stream transfer.

Mechanism note (pre-registered): in marginal mode the absent stream is
\emph{clamped} to explicit silence (never sampled), and under the
frozen adaptive decode the refinement loop exits every step ---
marginal-mode A.2 is therefore its AR pathway. Disclosed risk: an
all-silent partner is out-of-distribution for the melody/chord model
(chords are essentially always active in BOTH corpora of the combined
mix --- Nottingham tunes carry a rendered chord track throughout, and
the $\sim$13 melody-only tunes were excluded during preprocessing
precisely so the streams stayed paired --- so adding Nottingham does
NOT relieve this risk), while it is in-distribution for drum/non-drum,
where drum-less songs occur naturally. The two-task contrast is
therefore still the built-in diagnostic separating capacity
interference from distribution shift.

\subsection{E3 --- Conditional generation (horizon dose--response)}

The full ground-truth partner is given; the model generates the other
stream. The systems form an ordered conditioning-horizon spectrum
(Table~\ref{tab:e3-horizon}).

\begin{table}[ht]
\centering
\small
\begin{tabularx}{\textwidth}{@{}Lll@{}}
\toprule
Horizon of visible partner future & System & Streamable \\
\midrule
0 frames (same-frame) & A.2 conditional mode & yes \\
$k = 16$ frames (1 bar) & B.1 & yes, $k$-frame latency \\
unbounded (bidirectional) & C.1, C.2, Y & no \\
\bottomrule
\end{tabularx}
\caption{E3 conditioning-horizon spectrum.}
\label{tab:e3-horizon}
\end{table}

Scoring separates two notions of quality: \emph{fit to the given
stream} (primary --- what conditioning is for; computed by the metric
module on the pair $\langle$given ground truth, generated$\rangle$)
and agreement with the original target (secondary only --- many
accompaniments are valid for one partner). Hypotheses:
\begin{description}[nosep]
  \item[H-E3.1 (sanity gate, within-family).] Each system must beat
    its own co-generation output scored against the ground-truth
    partner --- same model, same decode, only the conditioning
    differs.
  \item[H-E3.2 (dose--response).] Monotone improvement along the
    horizon spectrum; the headline sub-question is the fraction of the
    $\mathrm{C.2}-\mathrm{A.2}$ gap B.1 recovers (a large fraction
    $\Rightarrow$ bounded, streamable conditioning suffices).
  \item[H-E3.3 (training vs.\ adapters).] Fully-trained duets vs.\
    adapters at fixed unbounded horizon --- deliberately registered
    without a directional prior, as both outcomes are informative.
\end{description}

\subsection{E4 --- Ablations}

\paragraph{E4a --- mechanism factorization.} A.2 differs from A.1 in a
coupled objective $+$ decode change; three contrasts separate them:
A.2($K{\ge}1$) vs.\ A.2($K{=}0$) isolates the decode-time refinement
(weights shared --- free); A.2($K{=}0$) vs.\ A.1 is the
training-objective parity check (expected parity; deficit $=$
objective tax, surplus $=$ beneficial auxiliary task); A.1 vs.\
A.2($K{=}4$) is the combined package (expected $\approx$ additive).
The refinement gain is predicted to be \emph{metric-selective}
(landing on H2, the fit block) --- a uniform gain would indicate
generic decode compute rather than the claimed mechanism. Disclosed
confound: A.1 and A.2 are separate training runs; their step counts
are recorded and any imbalance stated with its bias direction.

\paragraph{E4b --- decode schedule (melody/chord replication).}
$K \in \{0, 1, 4\}$ (expected: clearer monotone H2 benefit than on
drum/non-drum, since both streams are always active; a flat curve
would weaken the $K{=}4$ headline choice); adaptive early-exit on/off
at $K{=}4$ (expected null on melody/chord --- confirming the mechanism
is silence-specific); optional temperature-schedule comparison.

\paragraph{E4c --- MoE ablation.} A.2 vs.\ \sys{A.2-dense}
(compute-matched dense FFN), identical otherwise. Predicted deficit
profile is the mirror image of E4a's: landing on H3 (grammar) and
secondarily H1, with H2 least affected --- the coupling is carried by
the attention structure both variants share. Combined with the
observational expert-routing analysis, a positive result closes the
specialization story from both ends; a null is reported as
MoE-not-load-bearing.

\subsection{E5 --- Anticipation-horizon sweep (optional, gated)}

Only if E3 shows B.1 recovering a substantial fraction of the
unbounded-horizon gap: train B.1 at $k \in \{8, 16, 32\}$ and plot the
conditional fit against $k$, with A.1 as the $k{=}0$ point and C.2 as
the $k{=}\infty$ asymptote --- one dose--response figure locating the
knee of the lookahead curve.

\section{Evaluation methodology}

\subsection{Objective metrics (implemented: \code{eval\_metrics.py})}

All metrics are \emph{reference-calibrated}: distributional metrics
are Jensen--Shannon divergences against the ground-truth
continuation's distributions; fit metrics are deltas against the same
statistic computed on the reference pair; density is a ratio plus a
drift slope. Scoring is confined to the continuation window (frames
64--384). Reporting follows the pre-registered priority ordering
$\mathrm{H3} > \mathrm{H2} > \mathrm{H1}$, with one primary endpoint
per hypothesis; the remaining metrics are supporting diagnostics. The
module was validated on synthetic fixtures in which each engineered
failure mode (grammar chaos, stream death) trips exactly its own
hypothesis block.

\subsection{Listening test}

Blind pairwise A/B comparisons on three axes --- (a) overall
musicality, (b) inter-stream coherence, (c) 24-bar structure --- with
$\ge 3$ raters, randomized order, and samples pre-drawn at fixed seed
order (no curation). Win rates per axis with a sign test. The H4
prediction is about the \emph{profile} across axes, not the overall
win rate.

\subsection{Frozen decode settings}

One configuration per system, frozen before scoring: A.2 at $K{=}4$,
final temperature 0.9, top-$p$ 0.95, adaptive early-exit on ($K{=}4$
exercises the full trained refinement depth; the temperature schedule
is the prior listening-test selection; the adaptive guard covers the
known silent-frame regression and is itself ablated in E4b). All
other systems at temperature 1.0 with their script defaults. Any
change to these settings invalidates and reruns the affected grid.

\subsection{Statistics}

Paired per-song comparisons on identical prompts; three samples
averaged per song before testing; Wilcoxon signed-rank on per-song
means. With 10--15 evaluation songs, only large effects reach
conventional significance --- objective metrics are treated as
directional evidence and the listening test carries the headline
perceptual claims.

\section{Threats to validity (all pre-registered)}

\begin{enumerate}[nosep]
  \item \textbf{Genre mixing in the melody/chord training corpus} ---
    POP909 (pop piano) and Nottingham (folk tunes) differ in harmonic
    rhythm, voicing conventions and phrase structure, so the combined
    corpus does not represent one homogeneous ``melody/chord grammar.''
    This is accepted deliberately: it roughly doubles a corpus small
    enough that the first S1 finetune attempt overfit before its first
    validation measurement, and EVALUATION remains POP909-only so the
    reported numbers stay in the target domain. The cost is that H3
    (stream-appropriate grammar) is tested against a model whose
    training grammar is a genre mixture; a system could in principle
    lose H3 points for importing folk conventions into pop
    continuations. All melody/chord systems --- duet and baseline
    alike --- share the same mixture, so the COMPARISON stays matched
    even though the absolute H3 values are harder to interpret.
  \item \textbf{Y-mc domain position} --- the melody/chord YinYang is
    Nottingham-trained. It is reported as an external reference, never
    as a matched baseline. Note this threat CHANGED with the combined
    corpus: Nottingham is now part of every melody/chord system's
    training mix, so Y-mc is no longer the only system exposed to folk
    data --- but it is still the only one that never saw POP909, which
    is the evaluation domain. A POP909-matched YinYang finetune remains
    the identified optional upgrade.
  \item \textbf{E4a training-budget imbalance} --- A.1 and A.2 are
    separate runs; step counts are recorded and any imbalance
    disclosed with its bias direction.
  \item \textbf{E2 architecture non-constancy} --- there is no
    non-degenerate ``A.2 trained marginally,'' so E2 compares the
    joint system's marginal against the single-stream pipeline's best
    specialist; the specialist also sees half the tokens per song (an
    asymmetry that makes an A.2 win evidence of positive transfer,
    not a confound).
  \item \textbf{Multiple valid accompaniments} --- ground-truth
    agreement is never a primary conditional endpoint; fit-to-given
    is.
  \item \textbf{Small evaluation set} --- 10 POP909 $+$ 5--15 RWC
    songs bounds statistical power; mitigated by the paired design,
    three samples per cell, and the listening test as the perceptual
    arbiter.
  \item \textbf{Marginal-mode decode path} --- under the frozen
    adaptive setting, E2's A.2 outputs come from the AR pathway
    (refinement exits on the clamped-silent partner); stated so
    results are attributed to the right mechanism.
\end{enumerate}

\section{Execution roadmap}

\paragraph{Phase 0 --- prerequisites.} Data preparation for the
combined melody/chord corpus is complete: Nottingham was tokenized into
paired cp4 streams (melody = track 0, chords = track 1; ${\sim}13$
melody-only tunes excluded from BOTH streams so the pairing gate
passes), a program-tagged single-stream variant was built for the S*
baselines, and each was concatenated after POP909 --- POP909 FIRST and
unreordered, which preserves its songs' absolute indices and therefore
keeps the held-out evaluation ids held out under the index-mod-10 split.
\begin{enumerate}[nosep]
  \item A.2 melody/chord on the combined corpus --- \textbf{trained}.
  \item B.1 melody/chord on the combined corpus, $k{=}16$ ---
    \textbf{trained}.
  \item S1 on the combined merged/tagged corpus --- \textbf{trained}.
    \emph{Status note:} the first S1 attempt (20k steps, max-LR
    $2{\times}10^{-5}$, validation every 250 steps) overfit before its
    FIRST validation measurement --- ${\sim}20$ steps/epoch at this
    batch put that measurement ${\sim}12$ epochs in, past the minimum,
    and validation loss rose monotonically $0.37 \rightarrow 0.74$
    while train loss fell to $0.10$. The retuned recipe (2k-step
    schedule, max-LR $10^{-5}$, validation every 25 steps) captured the
    real dip on POP909-only: $0.39 \rightarrow {\sim}0.358$ over
    ${\sim}1$k steps, then a clean plateau. The same short recipe is
    used for the combined corpus and the specialists; because the
    combined corpus is roughly twice the size, its validation minimum
    is expected LATER in step count, so a still-descending curve at 2k
    steps calls for a fresh 4k-step run rather than a resume (the
    OneCycle LR has annealed to ${\sim}0$ by then).
  \item S-mel / S-chord specialists on the combined single-modality
    streams --- \textbf{to train}.
  \item E4c: A.2-dense on the combined corpus --- \textbf{to train}.
  \item Evaluation prompt folders for the ten held-out POP909 ids ---
    built automatically by \code{eval\_e1.sbatch}; RWC folders already
    exist.
  \item Locate the Y-dn / Y-mc checkpoints. (Metric and scoring
    modules: complete and dry-run verified.)
\end{enumerate}

\paragraph{Phase 1 --- generation sweeps.} All inference via existing
batch scripts: the duet lanes, the S*/specialist lanes, the YinYang
lanes, the E4a A.1 runs (existing checkpoint), and the E4b decode
grid. A.1-melchord training and the drum/non-drum specialists are
\emph{gated} on the corresponding first-task results.

\paragraph{Phase 2 --- scoring.} Manifest construction;
\code{eval\_metrics.py} over all outputs; per-mode comparison tables
in the pre-registered priority order.

\paragraph{Phase 3 --- listening test} on the E1/E3 shortlist; vote
collation.

\paragraph{Phase 4 --- write-up.} One results table per experiment
block; the E4b metric-vs-$K$ figure; the E4a factorization chart; the
E5 dose--response figure if gated in.

\medskip
\noindent Approximate additional training cost beyond runs already
underway: three short finetunes (S1, S-mel, S-chord $\approx$ a few
GPU-hours each) plus one full A.2-dense run ($\approx$ one
A.2-melchord budget). All evaluation sweeps are single-GPU inference
jobs.

\section{Planned deliverables}

\begin{itemize}[nosep]
  \item \textbf{T1}: E1 co-generation table (systems $\times$ H1--H3
    primaries $+$ listening win rates), per task.
  \item \textbf{T2}: E2 retention table (A.2 marginal vs.\
    specialist, margin verdict per stream).
  \item \textbf{T3}: E3 horizon table ordered by conditioning
    horizon, with the B.1 gap-recovery fraction highlighted.
  \item \textbf{F1}: E4a factorization bar chart (decode effect /
    objective effect / combined).
  \item \textbf{F2}: E4b metric-vs-$K$ curve.
  \item \textbf{F3}: expert routing by stream (observational
    companion to E4c).
  \item \textbf{T4}: E4c MoE table (A.2 vs.\ A.2-dense across
    hypothesis blocks).
  \item Listening-test protocol sheet and anonymized vote record.
\end{itemize}

\end{document}
